%==============================================================================
\chapter{Conclusiones y trabajo futuro}
\label{cap:conclusiones}
%==============================================================================

%------------------------------------------------------------------------------
\section{Conclusiones}
\label{sec:conclusiones}
%------------------------------------------------------------------------------

A lo largo de este trabajo se desarrolló un \textit{schedule} de ejecución
basado en portafolios para el JSSP clásico, que decide para cada instancia qué
configuraciones ejecutar, en qué orden y durante cuánto tiempo, a partir de
las características de la instancia y de modelos de aprendizaje automático.
Las conclusiones se organizan según los objetivos específicos y se refieren
al portafolio estudiado, formado por tres configuraciones de Chuffed y
Gecode, a las 236 instancias evaluadas y a un presupuesto de 300 segundos.

\paragraph{Objetivo 1: comparar enfoques de portafolio.} Se implementaron y
compararon, con la misma base de conocimiento y el mismo esquema de
validación, la selección de una configuración (\textit{Winner Takes All} con
vecinos más cercanos), el \textit{schedule} secuencial de SUNNY adaptado a
optimización, un \textit{schedule} estático y un \textit{schedule} construido
con aprendizaje automático. En el conjunto estudiado, seleccionar una sola
configuración da el mejor resultado en promedio, con un \textit{gap} medio de
13.12\,\% y el 87\,\% de la brecha entre el mejor \textit{solver} individual y
el VBS cerrada. Los \textit{schedules} quedan unas décimas por detrás en
promedio, pero son más robustos, puesto que el \textit{schedule} estático nunca
pierde más de 4.3 puntos de \textit{gap} frente a la mejor configuración,
mientras que \textit{Winner Takes All} llega a perder 7.8 cuando se equivoca. La
elección entre ambos enfoques depende, por tanto, de si se prioriza el
desempeño medio o la protección frente a errores de selección. La ventaja de
la selección se explica porque, en este conjunto, la mejor configuración es
fácil de predecir y las configuraciones siguen mejorando hasta el final del
presupuesto, con lo cual fraccionar el tiempo entre varias ejecuciones reduce el
margen de mejora de la configuración principal.

\paragraph{Objetivo 2: características de las instancias.} Se construyó un
conjunto de datos con 65 características genéricas de \texttt{mzn2feat} y 20
características propias del JSSP, sobre la carga de las máquinas, las
duraciones y la estructura de las rutas. Las estrategias que usan
características cierran entre el 77\,\% y el 87\,\% de la brecha, frente al
76\,\% de una regla que solo usa el tamaño de la instancia. Los dos grupos de
características dan resultados parecidos y su unión es la mejor opción para
\textit{Winner Takes All}. Las características propias, que se calculan en
menos de un segundo sin herramientas externas, constituyen una alternativa
práctica a \texttt{mzn2feat}.

\paragraph{Objetivo 3: portafolio con predicción y selección automática.} Se
implementó un prototipo funcional que, para una instancia nueva, calcula sus
cotas, extrae sus características, estima con tres modelos de bosques
aleatorios si cada configuración probará el óptimo, en cuánto tiempo y cuál
obtendrá la mejor solución, y construye con esas predicciones el orden y el
tiempo de ejecución. Los modelos son precisos (un AUC de 0.97 a 0.98 para
predecir la resolución y de 0.85 a 0.86 para identificar cuándo gana cada
una de las configuraciones de Chuffed), pero el \textit{schedule} construido con ellos resultó
equivalente en promedio al \textit{schedule} estático. En consecuencia, con
este portafolio, predecir bien el desempeño individual no se tradujo de forma
directa en un mejor reparto del tiempo.

\paragraph{Objetivo 4: validación experimental.} El portafolio se validó con
236 instancias clásicas de ocho familias, mediante validación
\textit{leave-one-out} y la simulación exacta de los \textit{schedules} a
partir de las trayectorias de los \textit{solvers}. Todas las estrategias de
portafolio superan al mejor \textit{solver} individual por entre 1.7 y 2.0
puntos de \textit{gap}, una diferencia estadísticamente significativa, y
recuperan todos o casi todos los óptimos que este pierde.
Como referencia externa, OR-Tools CP-SAT con el mismo modelo
obtiene un \textit{gap} medio de 3.97\,\%, muy por debajo de las estrategias
de este portafolio, lo que muestra que, en este caso, la selección entre estas
configuraciones mejora de forma consistente a sus miembros, pero no compensa
la distancia frente a CP-SAT.

\paragraph{Importancia de la formulación del modelo.} La mejora de mayor
magnitud en los experimentos provino de la formulación acotada del problema en
MiniZinc. Acotar el \textit{makespan} con una cota inferior y con la solución de
una heurística de despacho redujo el \textit{gap} medio de Chuffed de 648\,\% a
17\,\%, una ganancia muy superior a la que aporta cualquier estrategia de
selección sobre las mismas configuraciones. La comparación con los mismos
\textit{solvers} sin cotas (Tabla~\ref{tab:portafolios}) lo confirma: sin
cotas, ni eligiendo en cada instancia la mejor entre Chuffed y Gecode el
\textit{gap} medio baja del 24.5\,\%, y el margen entre el mejor
\textit{solver} y el VBS es de apenas 1.0 punto, mientras que con cotas los
enfoques de portafolio quedan cerca del 13\,\% con un margen de 2.24 puntos.
Asimismo, la heurística proporciona
una solución de respaldo que resguarda al portafolio frente a configuraciones
que no encuentran solución, y en seis instancias prueba la optimalidad sin
ejecutar ningún \textit{solver}. En este portafolio, aprovechar la
complementariedad entre sus miembros dependió de contar antes con una
formulación bien acotada para cada uno de ellos.

%------------------------------------------------------------------------------
\section{Limitaciones}
\label{sec:limitaciones}
%------------------------------------------------------------------------------

\begin{itemize}
  \item El portafolio activo se compone únicamente de dos \textit{solvers} de
    programación por restricciones en tres configuraciones y, aunque la
    heurística de despacho participa mediante el cálculo de las cotas y como
    solución de respaldo, no se incluyeron las metaheurísticas contempladas en
    el anteproyecto dentro del conjunto de algoritmos seleccionables.
  \item La evaluación se realizó sobre 236 instancias clásicas de la
    literatura, de modo que las conclusiones acerca de qué configuración
    conviene en cada caso están ligadas a esa distribución particular de tamaños
    y estructuras.
  \item El margen disponible entre el mejor \textit{solver} individual y el VBS
    es de 2.24 puntos de \textit{gap}, lo que acota la ganancia máxima que
    puede alcanzar cualquier estrategia de selección con estas configuraciones.
  \item Todas las conclusiones corresponden a un presupuesto fijo de 300
    segundos por instancia, por lo cual el reparto óptimo del tiempo y la
    conveniencia relativa de un \textit{schedule} podrían variar bajo otros
    límites de cómputo.
  \item Las configuraciones que integran un \textit{schedule} se ejecutan de
    manera independiente, sin transmitir la mejor solución encontrada de una
    etapa a la siguiente.
  \item CP-SAT se mantuvo fuera del portafolio como referencia externa con el
    propósito de estudiar la selección entre configuraciones de desempeño
    complementario.
\end{itemize}

%------------------------------------------------------------------------------
\section{Trabajo futuro}
\label{sec:trabajo_futuro}
%------------------------------------------------------------------------------

\begin{itemize}
  \item Transmitir la mejor solución encontrada por una configuración a la
    siguiente como nueva cota superior o como solución inicial, de manera que
    las etapas sucesivas del \textit{schedule} aprovechen el avance previo y el
    reparto del tiempo pueda superar a la selección de una sola configuración.
  \item Ampliar la composición del portafolio incorporando metaheurísticas como
    la búsqueda tabú, cuyo comportamiento en instancias grandes difiere del de
    los \textit{solvers} de programación por restricciones, junto con nuevas
    estrategias de búsqueda para los \textit{solvers} actuales.
  \item Emplear generadores de instancias que cubran regiones del espacio de
    características poco representadas en el conjunto clásico, especialmente
    aquellas donde la mejor configuración sea más difícil de predecir y un
    \textit{schedule} secuencial tenga mayor margen de aporte.
  \item Diseñar reglas de asignación de tiempo basadas en la predicción directa
    del \textit{gap} esperado en función del tiempo, en lugar de repartir el
    presupuesto en proporción a la probabilidad de que cada configuración sea la
    mejor.
  \item Extraer los descriptores de \texttt{mzn2feat} directamente sobre el
    modelo con cotas, de modo que las características reflejen el problema
    acotado que resuelven los \textit{solvers}.
  \item Extender el esquema de portafolio a variantes como el \textit{Flexible
    Job Shop} y a formulaciones multiobjetivo del problema, que quedaron fuera
    del alcance de este trabajo.
\end{itemize}
